\documentclass[12pt]{article}

\usepackage[a4paper, total={6in, 8in}]{geometry}
\usepackage{textcomp}

\begin{document}
\setlength{\parindent}{0pt}

\def \t {\textrightarrow}
\def \v {\vspace{1em}}
\def \bi {\begin{itemize}}
\def \ei {\end{itemize}}
\def \s[#1] {\section*{#1}}
\def \ss[#1] {\subsection*{#1}}
\def \sss[#1] {\subsubsection*{#1}}

\s[Montale]
\ss[I limoni]
È il testamento poetico dell'autore \t lirica da lettura del ruolo della poesia, e la funzione che la poesia ha nel 900 \\
Anche quasimodo riceve il premio nobel, e condivide la temperia culturale \\
Montale dice che non si impegna politicamente, ma contesta i poteri forti \t firma il manifesto degli intellettuali antifascisti \\
Il testo si caratterizza per alcune espressioni, e per un linguaggio semplice e piuttosto secco \t che fotografa la realta di quella liguria dove la natura brulla rimanda alla natura petrosa di Dante \\
Qua esplora il mondo della scelta poetica, e si esprime a riguardo dei poeti laureati

\v

"Ascoltami" \t eco di d annunzio nella pioggia del pineto \\
I limoni sono la ricchezza \t l eloquenza dei poeti laureati, il messaggio di d'annunizo rivistato e cambiato \t la scelta di montale che dichiara cosa si intenda per miracolo \\
Si schiude un portone \t e pochi versi prima (da 22): vedi questi silenzi \t quando siamo li li per capire il senso anche minimo della vita \t allora succede qualcosa che non permette di leggere oltre \\
Lo scoprire è uno sbagli di nautra \t essa ha un meccanicismo che imbriglia l uomo \t lo sbaglio di natura, l'anello della maglia che non tiene \\
Questo ci mette di fronte a una verita (non la verita) \t questa è una speranza, che potrebbe non essere rincompensata \\
La risorsa dell uomo ha le ali tarpate dalla natura \\
Versi dal 30 al 35 \t c'è delusione del nostro occhio, che cerca soluzione \t lo sguardo fruga \t la mente indaga, climax ascendente che porta a guardare + da vicino il reale, lo disunisce (immagine + forte del dividere, perche sottolinea mancanza di unione) \\
Nello smontarsi delle certezze, c e il conforto del profumo dei limoni \t la disturbata divinta si scontra con la divina indifferenza de "Il male di vivere ho incontrato" \\
Qua ci sono gia le anticipazioni

\v

Qui si trova una visione sistemica di quello che abbiamo fatto \\
L'immmagine dalla delusione, e la riflessione sui prodotti del progresso \\
La voce del 900, si interfaccia \t mostra come nulla delle provocazioni che lo hanno preceduto sono andate perse \\
Si va verso una visione empatica \t quando un giorno da un malchiuso portone \t parliamo di altro con tutte le immagine di limite viste, e parliamo di limite anche qua \\
Sensazioni uditive nell immagine che viene data alla fine del testo, trombe d'oro della solarita \t alle campane si sostituisce la solarità dei limoni \\
Sensazioni visive, gli alberi di una corte \t esprimono la spazialita, il contesto, una immagine ravvisaibilissima in liguria \\
Il cuore ha constatato che non ci sono formule che possano portare ad alemno una verita

\v

Poeti laureati \t sono i celeberrimi poeti, come d'annunzio \t e viene ripreso l incipit della pioggia del pineto \\
Insatura fa la parodia di d'annunzio \\
Complemento di limitazione \\
Nelle pozzznghere si trova l ilarita dei ragazzi che giocano \t vegetazione brulla ricorda ad alcuni paesaggi del Purgatorio (es. Manfredi) \\
Uccelli inghiottiti dall'azzurro \t che ascolta il sussurro dei rami \t complice di un anima che è stati \t siamo lontani dalla brezza quindi \t e abbiamo superato la lirica dell 800, tenendo fissi i presupposti esistenziali \t teologia e ontologia del negativo \\
Impregnano ogni dove \t ossimoro, una dolcezza infinita \\
I rami amici \t è esempio, anticipazione del tema ecologico \t Contini, gli alberi non sono alberi \\
Verso 18, divartite passioni per miracolo tace la guerra \t immagine della guerra aleggia, che arriva anche dove c era la pace che gli alberi dei limoni sapevano conferire \\
I poeti laureati \t a loro spettano rimore storiche, mentre a noi il miracolo al verso 18 \t le divertite passioni per miracolo tace la guerra \t cio che accade in questo luogo di sospensione, accade per miracolo \\
Il testo si puo leggere come un testo di intenzioni \\
Il problema letterari con d annunzio, quello esistenziale con il disfacimento dei valori \\
Espressione che ricorda svevo nel rompere tutto \t fetta di solarita, di luce mostrata \\
Lo sbaglio di natura, e tutto cio che accade per miracolo

\ss[La storia]
In "Dandoti il braccio" \t si trova nella sezione "Xenia" di Satura \t xenia erano i doni portati agli ospiti (da straniero, quindi zeno cosini) \\
Attenzione alla chiarezza del testo \t si ha superamento della visione manzoniana \\
FARE RIFERIMENTO A STORIA PER MATURITA \\
Demistificazione del meccanicismo \t impianto del vero poetico \\
L immagine della storia che viene patrocinata da questo testo, innette nella totale diffidenza nell insegnamento del passato \\
La sofferenza del poeta è incarnata negli oggetti \t e questo rende difficile da comprendere la sua poesia \t mentre in d annunzio si cerca eccesso di oggetti \\
Il dolore va attraversato, e la coscienza viene salvata \t il limite sovrasta comunque \t montale rappresenta nella lirica molto di quanto pirandello aveva mostrato nella narrativa

\s[Salvatore Quasimodo]
È il vero autentico ermetico tra gli altri poeti \t nasce nel 1901 \\
Salvatore quasimodo è appassionato di cultura straniera, letteratura e classicita (che traduce e studia) \\
Viene segnato da due eventi storici \t descrive le brutture della guerra \\
Sperimenta la testualita, e legge nel presente dell arco del 900 il peso della guerra sull intellettuale \t si interessa del tema della liberta dell intellettuale \\
Lui rappresenta infatti il fallimento dei poteri forti \\
Ha attaccamento alla sua terra, e si cela dietro questa la consapevolezza del 900 \t dice che bisogna rifare l uomo \t affermazione drammatica, ma veritiera \\
Uomo merita nuova lettura di se \t si legge infatti dalle poesie lo spirito del peso dello straniero, l assenza della libera \t che si lega alla condizione dell intellettuale nell impero (con seneca, tacito, quintilliano etc.) \\
Rifare l uomo, riconfigurare l uomo \t il suo ermetismo si legge nello sconcerto di fronte a un umanita dilaniata, e nella delusione dell uomo \\
Paesaggi e natura della sicilia devastata \t ma è l uomo che va ricostruito \\
Raccolte da mondadori, le + importanti:
\bi
  \item "Oboe sommerso" \t sezione "giorno dopo giorno"
  \item "Acque e terre"
  \item "Un uomo del mio tempo" = testamento poetico
\ei

\ss[Alle fronde dei salici]
Tratta da "oboe sommerso" individua il tema della liberta dell autore e dell intellettuale in generale \\
Il poeta si misura con qualcosa che lo sovrasta \t i testi che descrivono parte della sicilia e la sofferenza del vivere, e il dolore che si fa voce \\
Giuseppe De Robertis parla di una funzione di mediatore tra le esperienze della guerra, e la solarita della sicilia \t Elmengaldo parla della sofferenze ??

\v

Non solo le brutture della guerra, ma anche del progresso \t i figlio crocifissi dalla guerra, sono come cristo che si crocifigge \t immagine di salvezza rappresentata dai soldati \\
Sinestesie diverse, come urlo nero \t figlio crocifisso \t immagine della madre \t riprende foscolo \\
Lamento d'agnello dei fanciulli \t l'agnello sta per cristo \t i ragazzi si sacrificano in guerra, e richiamano il messaggio evangelico \\
Si legge quindi la crudelta della guerra \t climax ascendenti ultimi versi \\
Le fronde dei salici per voto, la poesia viene appesa \t le cetre vengono appesse \t la poesia puo ancora esistere? le cetre che ondeggiano simboleggiano come la poesia non puo essere un punto fermo \\
Non è la poesia che viene a meno, ma è schiacciata \\
Si coglie in questa un senso civile

\ss[Ed è subito sera]
La centralita dell uomo \t l uomo è solo sulla terra, e non c e consolazione \t trafitto da un raggio di solo \t questo non da conforto \\
Sono versi espressivi nella loro brevitas \t l unica frase di senso compiuto e quella del vero 1 \t trafitto è un participio \t si lega al resto come se fosse un attributo \\
Trafitto da un raggio di sole, ed e subito sera \\
Immagine di un uomo, ma chiunque \t non solo l intellettuale o il soldato \t umanita tutta, trafitta da un raggio di sole alla fine è sola \\
Niente + portare conforto all uomo, neanche il raggio di sole \t l illuminazione è il sole \\
Solitudine esistenziale, impossibilita di trovare il senso \t voce di non liberta, e l autore spiega come uomo di ogni tempo è lupo per gli altri \t evoca l immagine di Caino ed Abele, che non hanno esistato a usare la violenza uno contro l altro \\
È la natura dell uomo che ha lo spirito del combattere insito in se

\ss[Uomo del mio tempo]
È del 1946-1947 \\
Si rivolge a se stesso, all umanita, all uomo qualsiasi \\
Usa un linguaggio tecnico \t come per esempio la carlinga, che fa parte delle macchine da guerra \\
T'ho visto \t anafora \t eri tu \t senza amore e senza cristo, negazione di tutto \\
Eco fredda \t sinestesia, e anche nuvole di sangue (e anche visione allegorica che evoca la passione di cristo, quando il cielo si scuri) \\
Uccelli neri \t che bussano ai vetri in Carducci \\
Si legge qua dal primo fratricdio della bibbia \t caino uccide abele \\
Ricordare tecnicismo del linguaggio (pietra e fionda, preistoria), immutabilita, male dentro di noi \t si legge il testo in complementarieta con la "Storia" di Montale \\
Nasce nel 1901, muore nel 1968 \\
Riceve il nobel per la letteratura \t anche montale
\end{document}
